\documentclass[10pt, letterpaper]{article}

% Packages:
\usepackage[
    ignoreheadfoot,
    top=2 cm,
    bottom=2 cm,
    left=2 cm,
    right=2 cm,
    footskip=1.0 cm,
]{geometry}
\usepackage[explicit]{titlesec}
\usepackage{tabularx}
\usepackage{array}
\usepackage[dvipsnames]{xcolor}
\definecolor{primaryColor}{RGB}{0, 79, 144} % blue
\usepackage{enumitem}
\usepackage{fontawesome5}
\usepackage{amsmath}
\usepackage[
    pdftitle={Chunxu Han's CV},
    pdfauthor={Chunxu Han},
    pdfcreator={LaTeX with RenderCV},
    colorlinks=true,
    urlcolor=primaryColor
]{hyperref}
\usepackage[pscoord]{eso-pic}
\usepackage{calc}
\usepackage{bookmark}
\usepackage{lastpage}
\usepackage{changepage}
\usepackage{paracol}
\usepackage{ifthen}
\usepackage{needspace}
\usepackage{iftex}
\usepackage{multicol}
\usepackage{datetime2}
\usepackage{graphicx}

\ifPDFTeX
    \input{glyphtounicode}
    \pdfgentounicode=1
    \usepackage[T1]{fontenc}
    \usepackage[utf8]{inputenc}
    \usepackage{lmodern}
\fi

\usepackage[default, type1]{sourcesanspro}

\AtBeginEnvironment{adjustwidth}{\partopsep0pt}
\pagestyle{empty}
\setcounter{secnumdepth}{0}
\setlength{\parindent}{0pt}
\setlength{\topskip}{0pt}
\setlength{\columnsep}{0.15cm}
\makeatletter
\let\ps@customFooterStyle\ps@plain
\patchcmd{\ps@customFooterStyle}{\thepage}{
    \color{gray}\textit{\small Chunxu Han - Page \thepage{} of \pageref*{LastPage}}
}{}{}
\makeatother
\pagestyle{customFooterStyle}

\titleformat{\section}{
    \needspace{4\baselineskip}
    \Large\color{primaryColor}
}{}{}{
    \textbf{#1}\hspace{0.15cm}\titlerule[0.8pt]\hspace{-0.1cm}
}[]

\titlespacing{\section}{-1pt}{0.3 cm}{0.2 cm}

\newenvironment{highlights}{
    \begin{itemize}[
        topsep=0.10 cm,
        parsep=0.10 cm,
        partopsep=0pt,
        itemsep=0pt,
        leftmargin=0.4 cm + 10pt
    ]
}{
    \end{itemize}
}

\newenvironment{highlightsforbulletentries}{
    \begin{itemize}[
        topsep=0.10 cm,
        parsep=0.10 cm,
        partopsep=0pt,
        itemsep=0pt,
        leftmargin=10pt
    ]
}{
    \end{itemize}
}

\newenvironment{onecolentry}{
    \begin{adjustwidth}{0.2 cm + 0.00001 cm}{0.2 cm + 0.00001 cm}
}{
    \end{adjustwidth}
}

\newenvironment{twocolentry}[2][]{
    \onecolentry
    \def\secondColumn{#2}
    \setcolumnwidth{\fill, 4.5 cm}
    \begin{paracol}{2}
}{
    \switchcolumn \raggedleft \secondColumn
    \end{paracol}
    \endonecolentry
}

\newenvironment{threecolentry}[3][]{
    \onecolentry
    \def\thirdColumn{#3}
    \setcolumnwidth{1 cm, \fill, 4.5 cm}
    \begin{paracol}{3}
    {\raggedright #2} \switchcolumn
}{
    \switchcolumn \raggedleft \thirdColumn
    \end{paracol}
    \endonecolentry
}

\newenvironment{header}{
    \setlength{\topsep}{0pt}\par\kern\topsep\centering\color{primaryColor}\linespread{1.5}
}{
    \par\kern\topsep
}

\let\hrefWithoutArrow\href
\renewcommand{\href}[2]{\hrefWithoutArrow{#1}{\ifthenelse{\equal{#2}{}}{ }{#2 }\raisebox{.15ex}{\footnotesize \faExternalLink*}}}


\begin{document}
    \newcommand{\AND}{\unskip
        \cleaders\copy\ANDbox\hskip\wd\ANDbox
        \ignorespaces
    }
    \newsavebox\ANDbox
    \sbox\ANDbox{}

    % Portrait photo sits inline in the header, right side.
    % Copy portrait.jpg to the application folder before compiling.
    \begin{header}
        \begin{minipage}[t]{0.78\linewidth}
        \raggedright
        \fontsize{28 pt}{28 pt}
        \textbf{Chunxu Han}

        \vspace{0.3 cm}

        \normalsize
        \mbox{{\footnotesize\faMapMarker*}\hspace*{0.13cm}Peder Skrams Gade 16A, København K, Denmark}%
        \kern 0.25 cm%
        \AND%
        \kern 0.25 cm%
        \mbox{\hrefWithoutArrow{mailto:s220311@dtu.dk}{{\footnotesize\faEnvelope[regular]}\hspace*{0.13cm}s220311@dtu.dk}}%
        \kern 0.25 cm%
        \AND%
        \kern 0.25 cm%
        \mbox{\hrefWithoutArrow{tel:+45 52 73 06 85}{{\footnotesize\faPhone*}\hspace*{0.13cm}+45 52 73 06 85}}%
        \kern 0.25 cm%
        \AND%
        \kern 0.25 cm%
        \mbox{\hrefWithoutArrow{https://www.linkedin.com/in/chunxu-han/}{{\footnotesize\faLinkedinIn}\hspace*{0.13cm}chunxu-han}}%
        \kern 0.25 cm%
        \AND%
        \kern 0.25 cm%
        \mbox{\hrefWithoutArrow{https://github.com/chunxubioinfor}{{\footnotesize\faGithub}\hspace*{0.13cm}chunxubioinfor}}%
        \end{minipage}%
        \hfill%
        \begin{minipage}[t]{0.18\linewidth}
        \raggedleft
        \vspace{-0.2cm}
        \includegraphics[width=2.4cm, height=2.8cm, keepaspectratio=false]{portrait.jpg}
        \end{minipage}
    \end{header}

    \vspace{0.4 cm - 0.3 cm}

    \begin{onecolentry}
        \textit{I have experience in pharmaceutical consulting and quantitative research. After two years at DTU sharpening my analytical skills, I want to bring these skills back to consulting. I spoke with Baiyu Wang from your company and I am really interested in the job content. I hope you can consider my application and I am willing to share more.}
    \end{onecolentry}

    \vspace{0.3 cm}

    %% ===== TAILOR THIS: Relevant Experience (2–3 bold-lead paragraphs) =====
    %% Check source_materials/relevant_experience_bank.md first — reuse or adapt existing paragraphs.
    \section{Relevant Experience}

\begin{onecolentry}

\textbf{Pharmaceutical Consulting Experience} \\
During my internship at Illuminera (IQVIA), I worked on pharmaceutical consulting projects including market research, competitive landscape analysis, and social media listening. My academic background in biotechnology (BSc) and bioinformatics (MSc), combined with industry experience at Novo Nordisk, gives me a solid foundation for conducting market research within the bio and pharma domain. I understand both the scientific and business aspects of the pharmaceutical industry.

        \vspace{0.3 cm}

\textbf{Stakeholder \& KOL Engagement} \\
At Illuminera (IQVIA), I engaged with clients, stakeholders, and KOLs to gather insights and align project progress. I called KOLs from both hospitals and government institutions to gain insights for clients. I coordinated with multiple stakeholders and delivered strategic reports using PowerPoint. This experience taught me how to build professional relationships and extract valuable information from industry experts.

        \vspace{0.3 cm}

\textbf{Quantitative Research \& Data Analysis} \\
I have two years of experience in quantitative research at DTU HealthTech and Novo Nordisk. I contributed to developing an R package with a new statistical framework, conducted extensive quantitative modelling and benchmarking, and worked with large-scale healthcare datasets. At Novo, I developed statistical methods for assessing batch variability in cell therapy data. I work daily with R and Python and am comfortable designing and interpreting statistical analyses from scratch.

\end{onecolentry}

    %% ===== FIXED: Do not edit below this line =====

    \section{Education}

        \begin{twocolentry}{
    Lyngby, Denmark

    Sept 2023 -- March 2026
}
    \textbf{MSc in Bioinformatics and Systems Biology}, Technical University of Denmark
    \begin{highlights}
        \item I enrolled in both \textbf{Machine Learning} and \textbf{Deep Learning} courses, enabling me to develop an \textbf{AI-driven mindset} when formulating projects and developing new bioinformatics pipelines or tools.
    \end{highlights}
\end{twocolentry}

        \vspace{0.3 cm}

\begin{twocolentry}{
    Beijing, China

    Sept 2017 -- June 2021
}
    \textbf{BSc in Biotechnology}, Beijing Normal University
    \begin{highlights}
        \item \textbf{Coursework:} Built a solid foundation in \textbf{Biotech} with courses covering ecology, molecular biology, cell biology, and bioinformatics, alongside laboratory training.
    \end{highlights}
\end{twocolentry}


    \section{Work Experience}

        \begin{twocolentry}{
    Målev, Denmark

    Sep 2025 -- March 2026
}
    \textbf{Master Student}, Novo Nordisk

    \vspace{0.1cm}

    My master's thesis focuses on developing computational frameworks to evaluate batch-to-batch variability in single-cell RNA-seq data:

    \begin{highlights}

        \item \textbf{Methodology and Software Development:} Developed Batch Mixing Index (BMI), an entropy-based metric that quantifies per-cell batch mixing and incorporates \textbf{permutation-based statistical testing}, as well as scBatch, a \textbf{Python package} integrating seven complementary metrics for manufacturing batch variability assessment.

        \vspace{0.1cm}

        \item \textbf{AI/ML approaches:} Explored state-of-the-art data science approaches by applying the Deep Learning model to embed cell therapy batches into a shared latent space for variability assessment.

    \end{highlights}
\end{twocolentry}

        \vspace{0.3 cm}

        \begin{twocolentry}{
    Lyngby, Denmark

    Oct 2023 -- July 2025
}
    \textbf{Bioinformatics Research Assistant}, DTU HealthTech

    \vspace{0.1cm}

    I worked in Isoform Analysis Group for two years, which has been a transformative experience. I have grown from a student who only focuses on theoretical studies into a bioinformatician to deal with real-world problems:

    \begin{highlights}
        \item \textbf{Development and Maintenance of Tools:} Contributed to developing and updating the \href{https://github.com/kvittingseerup/IsoformSwitchAnalyzeR}{IsoformSwitchAnalyzeR} R package. This experience sharpened my coding skills and can-do mindset, especially when handling issues with tools I'd never used before. I pushed myself to attempt and now I'm comfortable with the steep learning curve.

        \vspace{0.1cm}

        \item \textbf{Bioinformatics Pipelines:} Developed and optimized bioinformatics pipelines, supporting others in running analyses. I mastered multi-tasking, learned to leverage \textbf{HPC} for parallel processing, and \textbf{Git} to achieve version control.

        \vspace{0.1cm}

        \item \textbf{Ad-hoc Work:} Managed various ad-hoc tasks, from literature reviews to data organization, learning how to \textbf{prioritize and organize} work effectively in a fast-paced research environment.
    \end{highlights}
\end{twocolentry}

        \vspace{0.3 cm}

        \begin{twocolentry}{
    Shanghai, China

    Feb 2023 -- May 2023
}
    \textbf{Consulting Intern}, Illuminera, IQVIA

    \vspace{0.1cm}

    My first experience in a business setting, where I gained insights into the pharmaceutical market and developed professional skills:

    \begin{highlights}
        \item \textbf{Qualitative and Quantitative Research:} Conducted market trend and competitor analysis to formulate actionable strategies to clients. I also gained a deeper understanding of the pharmaceutical market in China and globally.
        

        \vspace{0.1cm}

        \item \textbf{Social Media Listening:} Led a social media listening project using platforms such as RedNote and online hospital communities, collecting and analysing large-scale user data to identify product gaps, map patient and consumer preferences, and generate strategic recommendations for \textbf{product positioning} and \textbf{targeted commercial campaigns}.

        \vspace{0.1cm}

        \item \textbf{Client and KOLs Engagement:} Built professional relationships with clients and KOLs and coordinated with stakeholders to align project progress and delivery.

        \vspace{0.1cm}

        \item \textbf{Data-Driven Deliverables:} Utilized tools like PowerPoint and Power BI to deliver strategic reports, transforming complex data into actionable business insights. I learned how to formulate the business problems and storytelling to the clients.
    \end{highlights}
\end{twocolentry}

        \vspace{0.3 cm}


    %% Publications section commented out for data/tech roles.
    %% Uncomment if the role values research output.
%     \section{Publications}
%
%         \begin{onecolentry}
%             \begin{itemize}
%                 \item \textbf{Han, C.}, Gilis, J., Iriondo Delgado, E., Clement, L., \& Vitting-Seerup, K. \textit{IsoformSwitchAnalyzeR v2: Analysis of Functional Isoform Changes in Long-read and Single-cell Sequencing Data}. bioRxiv (2025). \href{https://doi.org/10.64898/2025.12.08.693027}{doi:10.64898/2025.12.08.693027}
%
%                 \vspace{0.1 cm}
%
%                 \item Mikkelsen, J. M., \textbf{Han, C.}, Kanakoglou, D. S., Tiberi, S., \& Vitting-Seerup, K. \textit{Beyond the Gene in Genetics: How Isoform-Resolved Analysis Empowers the Study of Both Common and Rare Genetic Variation}. medRxiv (2025). \href{https://doi.org/10.1101/2025.03.31.25324931}{doi:10.1101/2025.03.31.25324931}
%             \end{itemize}
%         \end{onecolentry}


    \section{Skills and Hobbies}

        \begin{onecolentry}
            \textbf{Programming:} Python, R, Bash, SQL, Snakemake
        \end{onecolentry}

        \vspace{0.2 cm}

        \begin{onecolentry}
            \textbf{Languages:} English (proficient), Chinese (native), Danish (Basic)
        \end{onecolentry}

        \vspace{0.2 cm}

        \begin{onecolentry}
            \textbf{Soft Skills:} Communication, Teamwork, Client Engagement, Multitasking, Critical Thinking, Strong Work Ethic
        \end{onecolentry}

        \vspace{0.2 cm}

        \begin{onecolentry}
            \textbf{Hobbies:}
            I've been playing \textbf{badminton} since an early age. It's a hobby that keeps me focused and energized. During holidays, I enjoy \textbf{traveling} around Europe to experience new food and cultures. I also love \textbf{hiking}, both for the chance to connect with nature and for the challenge of seeing hidden scenery that isn't easily accessible. After moving to Denmark, I discovered a new passion: \textbf{sailing}. Last summer, we sailed to Ven and Hallands Väderö, and it was a completely new experience and I feel free to ride the infinite waves.
        \end{onecolentry}


\end{document}
